\chapter{Envelope and extremizer}\label{chap:envelope}

This chapter explains why the envelope from Chapter~\ref{chap:stoploss} is
sharp. The constants $a$ and $p_0$ are tuned so that the piecewise formula for
$J_G$ is tangent to the Gaussian comparator at the correct contact point, and a
canonical probability measure $\mu^\star$ is then built to saturate the tail cap
and the envelope exactly.

\section{The envelope parameters}

\begin{lemma}[The threshold parameter \(a\) and the tangent line]\label{lem:a-parameter}
\lean{exists_unique_a}\leanok
\lean{J_ge_tangent_line}\leanok
There is a unique $a > t_0$ such that $H(a)=B$, and the resulting envelope
$J_G$ lies above the affine function $u \mapsto B - p_0 u$ on $[0,\infty)$.
\end{lemma}

\begin{proof}\leanok
The map $H(x)=x s_G(x)+\int_x^\infty s_G(t)\,dt$ is strictly decreasing on
$(t_0,\infty)$ and decays to $0$ at infinity, while $H(t_0)=A=2B$.
This gives a unique solution to $H(a)=B$. The tangent-line lower bound for
$J_G$ is then a direct consequence of the piecewise definition of the envelope:
on $[0,a]$ the envelope is exactly the affine line, and on $[a,\infty)$ the
tail-integral branch stays above its common value at $u=a$.
\end{proof}

\section{The canonical extremal measure}

\begin{theorem}[Extremal measure saturating the tail cap]\label{thm:extremizer-tail}
\lean{exists_extremal_measure}\leanok
\lean{muStar_isProbability}\leanok
\lean{mean_muStar}\leanok
\lean{absTail_muStar_eq_s}\leanok
Lean constructs an explicit probability measure $\mu^\star$ with mean zero such
that
\[
  \Prb_{\mu^\star}(|x|>t) = s_G(t) \qquad (t \ge 0).
\]
This is the canonical witness showing that the stop-loss envelope is sharp.
\end{theorem}

\begin{proof}\leanok
The formalization defines a piecewise density whose right tail matches the
sub-Gaussian cap and whose left side is adjusted so that the total mass is one
and the mean vanishes. The key calculations show that the right and left tail
integrals match the desired formulas exactly, after which the probability and
mean-zero statements follow by direct integration.
\end{proof}

\begin{theorem}[The extremizer attains the envelope]\label{thm:extremizer-stoploss}
\lean{stopLoss_muStar_eq_J}\leanok
\uses{thm:envelope,thm:extremizer-tail}
For every $u \ge 0$,
\[
  \mathrm{stopLoss}_{\mu^\star}(u) = J_G(u).
\]
\end{theorem}

\begin{proof}\leanok
Once the tail identity from Theorem~\ref{thm:extremizer-tail} is known, the
layer-cake formula turns the stop-loss of $\mu^\star$ into the integral of the
tail cap itself. That integral is exactly the definition of the envelope on the
tail branch, while the affine branch is forced by the centering condition. This
shows that the upper envelope from Chapter~\ref{chap:stoploss} is not merely an
upper bound but the true optimum.
\end{proof}
